\section{Discussion and Limitations}

The results distinguish an internal architectural improvement from a
competitive compression system. Original-grid support improves this
architecture relative to its own unconstrained router, and a three-seed
CIFAR campaign finds the same sign against global, flattened, and
degree-matched controls. Yet the full ImageNet model loses accuracy to the
selected dense baseline at 224, 384 and 512 pixels, its 224-pixel
implementation is slower in guarded synthetic tests, and post-training ToMe
provides a stronger accuracy baseline at every matched budget in the sweep.
Optimizer-side 150-epoch probes do not close the selected-configuration gap.
These findings motivate revisiting the computation path, the compression
schedule, and representation initialization; they do not show that learning a
merge criterion is intrinsically inferior.

The evidence is limited in several ways. All ImageNet training runs use seed
42, and the validation set serves both checkpoint selection and reporting.
We have no across-seed ImageNet uncertainty estimate or ImageNet per-image
predictions for paired significance tests. Section~\ref{sec:cifar-seeds} adds
three-seed training on resized CIFAR-100, which is not native-resolution
ImageNet.
Geometry and candidate cardinality remain coupled on ImageNet. The
300-epoch \mn{} configuration benefits from an explicit learning-rate screen,
while the dense baseline is not subjected to an equally broad screen. The
$\lambda=4$ training variant changes multiple architectural parameters, and
test-time budget changes do not replace a clean training-budget ablation.
The 512-pixel \mn{} run is the radius-scaled branch seeded from the
radius-5.143 384-pixel EMA file, not a radius-3 continuation.

The experiment artifacts also bound reproducibility. Raw metrics, resolved
arguments, launch records, environment information, benchmark scripts and
checkpoint receipts are available in the campaign packet. Company policy
prevents transfer of the trained ImageNet weights, and a release commit alone
does not capture every change in the collaborator's working tree. The companion
analysis preserves the source identities and distinguishes recorded results
from independent follow-up diagnostics. The training environment uses PyTorch
2.7.1 with CUDA 11.8, timm 0.9.11 and FlashAttention 2.7.3, which differs from
the local release's pinned environment. CPU profiles and H20 synthetic timings
cannot establish performance on every accelerator or under deployment load.
The local follow-up resolves the ToMe attention-policy and patch-count
mismatch for synthetic architecture timing. A verified trained-checkpoint
joint frontier remains unavailable.

Additional initialization seeds, a matched-final-budget progressive schedule,
and Base-scale training remain useful follow-up directions, but they are not
required to support the present method claim. The completed single-seed
diagnostics establish initialization sensitivity and no observed gain from one
more aggressive latent schedule; they do not establish equal-compute or
matched-budget conclusions. The submission depends only on two unavailable
checkpoint-side outputs: a common-protocol numeric packet for the main table
and component interventions, and passive routing traces for the final
visualization.
